% Canonical Companion Part IV source figure
% Companion problem: CP-IV-0089
% Source problem: IMP-DL-E88547CF67-P030
% Source: 3611-L4212
% This file preserves the source TikZ/figure block verbatim except for this provenance header.
\begin{figure}[H]
		\centering
		\begin{tikzpicture}[>=Latex, line cap=round, line join=round, scale=1.05]
			
			% Colors matching the previous chart
			\definecolor{mapblue}{RGB}{36,99,181}      % blue for maps/arrows
			\definecolor{datagreen}{RGB}{34,139,34}    % green for g-value
			\definecolor{phasered}{RGB}{176,54,54}     % (not essential here)
			\definecolor{sphgray}{RGB}{245,245,245}
			
			% Sphere
			\fill[sphgray] (0,0) circle (1.8);
			\draw[thick] (0,0) circle (1.8);
			\node at (0,2.35) {\small $\mathbb S^2$};
			
			% North pole
			\fill[mapblue] (0,1.8) circle (1.8pt);
			\node[above] at (0,1.8) {\small $N_0$};
			
			% Equatorial plane (identified with C)
			\draw[thick] (-3.2,-1.9) -- (3.2,-1.9);
			\node[below right] at (3.2,-1.9) {\small $\mathbb C$ (stereographic plane)};
			
			% A point on the sphere = normal vector N(p)
			\fill[mapblue] (1.0,0.9) circle (1.8pt);
			\node[above right] at (1.0,0.9) {\small $N(p)$};
			
			% Stereographic line from N0 through N(p) to plane
			\draw[->, very thick, mapblue] (0,1.8) -- (1.9,-1.9);
			
			% Image point g(p) in complex plane
			\fill[datagreen] (1.9,-1.9) circle (1.8pt);
			\node[below right] at (1.9,-1.9) {\small $g(p)=\sigma(N(p))$};
			
			% Labels for maps
			\node[align=left] at (-3.3,0.9) {\small \textcolor{mapblue}{$N:\text{(surface)}\to\mathbb S^2$}\\[-1pt]
				\small \textcolor{mapblue}{Gauss map}};
			\node[align=left] at (-3.6,0.1) {\small \textcolor{mapblue}{$\sigma:\mathbb S^2\setminus\{N_0\}\to\mathbb C$}\\[-1pt]
				\small \textcolor{mapblue}{stereographic projection}};
			\node[align=left] at (-3.15,-0.85) {\small \textcolor{datagreen}{$g=\sigma\circ N$}\\[-1pt]
				\small \textcolor{datagreen}{normal direction as a complex number}};
			
		\end{tikzpicture}
		\caption{The Weierstrass function $g$ is the Gauss map written in complex coordinates: $g=\sigma\circ N$.}
	\end{figure}
